\documentclass[11pt]{article}

\usepackage[margin=1.08in]{geometry}
\usepackage[T1]{fontenc}
\usepackage{lmodern}
\IfFileExists{microtype.sty}{\usepackage{microtype}}{}
\usepackage{amsmath,amssymb,amsthm,mathtools}
\usepackage{booktabs}
\usepackage{enumitem}
\usepackage{xcolor}
\usepackage{hyperref}
\hypersetup{
  colorlinks=true,
  linkcolor=blue!55!black,
  citecolor=blue!55!black,
  urlcolor=blue!55!black,
  pdftitle={Sparse Multi-Base Representations and Asymptotic Existence of Hadamard Matrices}
}

\newtheorem{theorem}{Theorem}[section]
\newtheorem{proposition}[theorem]{Proposition}
\newtheorem{corollary}[theorem]{Corollary}
\newtheorem{lemma}[theorem]{Lemma}
\newtheorem{remark}[theorem]{Remark}
\newtheorem{definition}[theorem]{Definition}

\newcommand{\Hcal}{\mathcal{H}}
\newcommand{\N}{\mathbb{N}}
\newcommand{\R}{\mathbb{R}}
\newcommand{\eps}{\varepsilon}
\newcommand{\wt}{\operatorname{wt}}
\newcommand{\ee}{\mathrm{e}}

\title{\textbf{Sparse Multi-Base Representations and\\Asymptotic Existence of Hadamard Matrices}}
\author{Joshua Mark\thanks{ORCID: \href{https://orcid.org/0009-0003-5905-2382}{0009-0003-5905-2382}. AI assistance in formalization and preparation is disclosed in the accompanying metadata.}}
\date{2 November 2026 --- formalization-aligned revision}

\begin{document}
\maketitle

\begin{abstract}
For an odd positive integer $p$, let
\[
\tau(p):=\min\{t\ge 0: \text{a Hadamard matrix of order }2^t p\text{ exists}\}.
\]
We combine a theorem of Krenn, Suppakitpaisarn and Wagner on sparse multi-base representations with Craigen's complementary-sequence construction for Hadamard matrices.  The resulting transfer gives the uniform bound
\[
\tau(p)=O\!\left(\frac{\log p}{\log\log p}\right),
\]
which is asymptotically smaller than every bound of the form $c\log p+O(1)$.  The key observation is that every integer is a sum of $O(\!\log n/\log\log n)$ distinct numbers $2^a10^b$, while every $2^a10^b$ is the length of a binary Golay complementary pair.

We also address an additive-number-theory problem posed by Livinskyi.  If $S(b_1,\ldots,b_r)$ is the multiplicative monoid generated by fixed multiplicatively independent integers $b_1=2,b_2,\ldots,b_r$, and $\ell_m(S)$ is the smallest positive integer not expressible as a sum of at most $m$ elements of $S$, then
\[
\exp(c m\log m)\le \ell_m(S)\le
\exp\bigl((r-1+o(1))m\log m\bigr)
\]
for some $c>0$.  In particular,
$\ell_m(S(2,10,26))=\exp(\Theta(m\log m))$.
As consequences of the Hadamard construction we obtain the quantitative gap estimate
\[
\gamma(x)\le \exp\!\left(O\!\left(\frac{\log x}{\log\log x}\right)\right)=x^{o(1)}
\]
for consecutive Hadamard orders.  Combining this with the 2025 argument of Alexeev, Jasper and Mixon yields, for every $\delta>0$, the asymptotic approximate-Hadamard bound
\[
\kappa(n)\le 1+n^{-1/4+\delta}
\]
for all sufficiently large $n$.
\end{abstract}

\noindent\fbox{\parbox{0.94\textwidth}{\small
\textbf{Status and priority note.}
The mathematical deductions in this draft are presented for verification and development.  A targeted literature search located the ingredients below but did not locate a paper combining the 2020 multi-base theorem with the Craigen--Livinskyi Hadamard construction in this way.  This is \emph{not} an assertion of priority.  Before submission, the novelty claim should be checked through MathSciNet, zbMATH, Google Scholar citation chains, and, ideally, correspondence with authors working on asymptotic Hadamard existence.  A qualitative subpolynomial Hadamard-gap result was already attributed in 2013 to an unpublished manuscript of W.~D.~Smith; see Section~\ref{sec:gaps}.
}}

\section{Introduction}
A Hadamard matrix of order $n$ is a matrix $H\in\{\pm1\}^{n\times n}$ satisfying
\[
HH^{\mathsf T}=nI_n.
\]
Except for orders $1$ and $2$, divisibility by $4$ is necessary.  The Hadamard conjecture asserts that this necessary condition is sufficient: a Hadamard matrix should exist at every order $4k$.

A classical asymptotic relaxation fixes an odd integer $p$ and asks how many powers of two must be adjoined before a Hadamard order is guaranteed.  Wallis proved the first general logarithmic bound in 1976~\cite{Wallis1976}; Craigen improved the coefficient using signed groups and complementary sequences~\cite{Craigen1995}.  More recently, Du and Jiang proved that for every odd $m$ there is a Hadamard matrix of order $2^t m$ with
\begin{equation}\label{eq:dujiang}
 t=6\left\lfloor \frac{1}{40}\log_2 m\right\rfloor+10
   =\frac{3}{20}\log_2 m+O(1).
\end{equation}
Their paper also points out that one of the stronger numerical bounds stated in Livinskyi's 2012 thesis relies on an earlier theorem that is incorrect~\cite{DuJiang2024}.  The argument developed here does not use that result: it uses only the binary-Golay transfer theorem originating with Craigen and quoted as Theorem~3.10 in Livinskyi's thesis~\cite{Craigen1995,Livinskyi2012}.

Livinskyi also proposed a different route.  Write
\[
S(n_1,\ldots,n_k)=\{n_1^{a_1}\cdots n_k^{a_k}:a_i\in\N_0\}
\]
and let $\ell_m(S)$ be the smallest positive integer not representable as a sum of at most $m$ elements of $S$.  Since
\[
S(2,10,26)
\]
consists of binary Golay lengths, Livinskyi observed that strong lower bounds on $\ell_m(S(2,10,26))$ would translate into stronger bounds for $\tau(p)$, but no asymptotic estimate for this covering function was obtained there~\cite{Livinskyi2012}.

The missing ingredient appears in a different part of number theory.  Krenn, Suppakitpaisarn and Wagner proved that, for any fixed set of at least two multiplicatively independent bases satisfying a mild representability assumption, the worst-case minimal Hamming weight of a multi-base representation is
\[
\Theta\!\left(\frac{\log n}{\log\log n}\right).
\]
Their upper bound is realized by a greedy algorithm and may be chosen with digits only $0$ and $1$~\cite{Krenn2020}.  Applying their theorem to the bases $2$ and $10$ produces sparse sums of binary Golay lengths, exactly the input required by Craigen's Hadamard construction.

The main contributions of this note are:
\begin{enumerate}[leftmargin=2em]
\item a sparse Golay decomposition of every positive integer using $O(\log n/\log\log n)$ binary Golay lengths;
\item the uniform Hadamard-existence estimate $\tau(p)=O(\log p/\log\log p)$ for all odd $p$;
\item the order of growth $\log \ell_m(S(2,10,26))=\Theta(m\log m)$, answering Livinskyi's proposed additive-number-theory direction at the level of order of magnitude;
\item a quantitative subpolynomial bound on gaps between Hadamard orders; and
\item an unconditional improvement, via Alexeev--Jasper--Mixon~\cite{Alexeev2025}, of their approximate-Hadamard exponent from every $\alpha<17/92$ to every $\alpha<1/4$.
\end{enumerate}

\section{Golay pairs, complementary sequences, and multi-base representations}

\subsection{Binary Golay lengths}
For a binary sequence $A=(a_0,\ldots,a_{n-1})\in\{\pm1\}^n$, write its aperiodic autocorrelation at positive shift $s$ as
\[
N_A(s)=\sum_{j=0}^{n-s-1}a_j a_{j+s}.
\]
A pair $(A,B)$ of binary sequences of the same length $n$ is a \emph{Golay complementary pair} if
\[
N_A(s)+N_B(s)=0\qquad(1\le s<n).
\]
The corresponding $n$ is called a binary Golay length.  Golay introduced complementary series~\cite{Golay1961}; Turyn's multiplication construction implies the following standard family~\cite{Turyn1974}.

\begin{proposition}[Golay--Turyn]\label{prop:golaylengths}
For every $a,b,c\in\N_0$, there exists a binary Golay complementary pair of length
\[
2^a10^b26^c.
\]
The length $1$ is understood as the trivial pair.
\end{proposition}

Only the subfamily $2^a10^b$ will be needed for the main theorem.  The larger monoid $S(2,10,26)$ is useful for the covering problem in Section~\ref{sec:covering}.

\subsection{Craigen's transfer to Hadamard matrices}
The following construction is due to Craigen~\cite{Craigen1995}; the convenient formulation below is Theorem~3.10 of Livinskyi~\cite{Livinskyi2012}.

\begin{theorem}[Craigen]\label{thm:craigen-transfer}
Suppose that
\[
A_1,B_1,\ldots,A_r,B_r
\]
are $2r$ binary sequences whose total aperiodic autocorrelation is zero, with $A_i$ and $B_i$ both of length $l_i$.  Put $L=\sum_{i=1}^r l_i$.  Then there exists a Hadamard matrix of order
\[
4^{r+1}(2L+1).
\]
\end{theorem}

If each $(A_i,B_i)$ is a Golay pair, the hypothesis is automatic because the autocorrelation sums vanish pairwise.

\subsection{Sparse multi-base representations}
Fix integers $b_1,\ldots,b_s>1$.  Their power-product monoid is
\[
\mathcal B=\{b_1^{a_1}\cdots b_s^{a_s}:a_i\in\N_0\}.
\]
For a finite digit set $D\subset\N_0$, a multi-base representation of $n$ is
\[
 n=\sum_{B\in\mathcal B} d_BB,
 \qquad d_B\in D,
\]
with only finitely many nonzero coefficients.  Its Hamming weight is the number of nonzero $d_B$.  The bases are \emph{multiplicatively independent} if
\[
b_1^{z_1}\cdots b_s^{z_s}=1,\qquad z_i\in\mathbb Z,
\]
implies $z_1=\cdots=z_s=0$.

We use the following result of Krenn, Suppakitpaisarn and Wagner~\cite{Krenn2020}.

\begin{theorem}[Krenn--Suppakitpaisarn--Wagner]\label{thm:ksw}
Let $s\ge2$.  Suppose $b_1,\ldots,b_s$ are multiplicatively independent and the finite nonnegative digit set $D$, containing $0$, permits a multi-base representation of every positive integer.  There are constants $K_1,K_2>0$ such that:
\begin{enumerate}[label=(\roman*),leftmargin=2em]
\item every $n>2$ has a representation of Hamming weight at most
\[
K_1\frac{\log n}{\log\log n};
\]
\item for infinitely many $n$, every such representation has Hamming weight at least
\[
K_2\frac{\log n}{\log\log n}.
\]
\end{enumerate}
In the applications below we take $D=\{0,1\}$ and verify its representability hypothesis separately; no conversion of an arbitrary digit set into $\{0,1\}$ is asserted.
\end{theorem}

The constants depend on the bases and digits.  The proof is asymptotic; in particular, the preprocessing threshold in the greedy algorithm may be very large~\cite{Krenn2020}.

\section{Sparse decompositions into binary Golay lengths}
The crucial specialization of Theorem~\ref{thm:ksw} is unexpectedly simple.

\begin{lemma}\label{lem:independence}
The integers $2$ and $10$ are multiplicatively independent.
\end{lemma}
\begin{proof}
If $2^a10^b=1$ for $a,b\in\mathbb Z$, then
\[
2^{a+b}5^b=1.
\]
Unique factorization forces $b=0$, and then $a=0$.
\end{proof}

\begin{theorem}[Sparse Golay decomposition]\label{thm:sparse-golay}
There is a constant $K>0$ such that every integer $q>2$ can be written
\begin{equation}\label{eq:sparse-golay}
q=\sum_{j=1}^{r}2^{a_j}10^{b_j},
\qquad a_j,b_j\in\N_0,
\end{equation}
with the summands distinct and
\[
r\le K\frac{\log q}{\log\log q}.
\]
Every summand in \eqref{eq:sparse-golay} is a binary Golay length.
\end{theorem}

\begin{proof}
Apply Theorem~\ref{thm:ksw} with bases $(2,10)$ and digit set $D=\{0,1\}$.  Multiplicative independence follows from Lemma~\ref{lem:independence}.  The representability hypothesis holds because ordinary binary expansion is already a multi-base representation using only the power-products $2^a10^0$.  The theorem therefore gives a representation with $O(\log q/\log\log q)$ nonzero $0/1$ digits, which is exactly \eqref{eq:sparse-golay}.  Proposition~\ref{prop:golaylengths} gives a binary Golay pair at each such length.
\end{proof}

\begin{remark}\label{rem:morebases}
The same argument works with any fixed collection of multiplicatively independent binary Golay generators containing $2$.  The pair $(2,10)$ is enough for the asymptotic order proved here and keeps the transfer transparent.
\end{remark}

\section{A sublogarithmic exponent for asymptotic Hadamard existence}
Let $\Hcal$ denote the set of Hadamard orders.  For odd $p$ define
\[
\tau(p)=\min\{t\ge0:2^tp\in\Hcal\}.
\]

\begin{theorem}[Main Hadamard bound]\label{thm:main}
There is a constant $C>0$ such that every sufficiently large odd positive integer $p$ satisfies
\begin{equation}\label{eq:main}
\tau(p)\le C\frac{\log p}{\log\log p}.
\end{equation}
Equivalently, every sufficiently large odd $p$ divides a Hadamard order
\[
p\,\exp\!\left(O\!\left(\frac{\log p}{\log\log p}\right)\right)=p^{1+o(1)}.
\]
\end{theorem}

\begin{proof}
Put
\[
q=\frac{p-1}{2}.
\]
By Theorem~\ref{thm:sparse-golay}, write
\[
q=l_1+\cdots+l_r,
\qquad
r=O\!\left(\frac{\log q}{\log\log q}\right),
\]
where every $l_i$ is a binary Golay length.  Choose a Golay pair $(A_i,B_i)$ of length $l_i$ for each $i$.  The $2r$ sequences have zero total aperiodic autocorrelation, so Theorem~\ref{thm:craigen-transfer} gives a Hadamard matrix of order
\[
4^{r+1}(2q+1)=4^{r+1}p=2^{2r+2}p.
\]
Hence
\[
\tau(p)\le 2r+2
=O\!\left(\frac{\log p}{\log\log p}\right),
\]
which is \eqref{eq:main} after enlarging the constant.
\end{proof}

\begin{remark}[Comparison with the recent logarithmic bound]
Du and Jiang's formula \eqref{eq:dujiang} is explicit and is likely far stronger at computationally accessible sizes.  Theorem~\ref{thm:main} is instead an asymptotic statement: since
\[
\frac{\log p}{\log\log p}=o(\log p),
\]
it eventually improves every positive constant multiple of $\log p$.  No useful crossover point follows from the present proof because the constants inherited from Theorem~\ref{thm:ksw} are not explicit in a practical form.
\end{remark}

\begin{remark}[A limitation of the fixed multi-base mechanism]\label{rem:optimality}
The lower-bound half of Theorem~\ref{thm:ksw} shows that, for the fixed bases $(2,10)$ and digits $\{0,1\}$, infinitely many integers require
$\Omega(\log n/\log\log n)$ distinct power-products.  Allowing repeated summands does not circumvent this order: whenever two equal elements of $S(2,10)$ occur, they may be merged into their double, which remains in $S(2,10)$; iterating produces a representation by distinct elements without increasing the number of terms.  Thus the pair count in Theorem~\ref{thm:sparse-golay} has the correct worst-case order within this fixed-monoid representation framework.  A further asymptotic improvement would require additional structure---for example, a richer family of Golay lengths, a different additive representation theorem, or a Hadamard transfer whose power-of-two overhead grows more slowly than linearly in the number of complementary pairs.
\end{remark}

\section{Livinskyi's covering function}\label{sec:covering}
Livinskyi defined, for a set $S\subset\N$,
\begin{equation}\label{eq:ell-def}
\ell_m(S):=\min\{n\ge1:n\text{ is not a sum of at most }m\text{ elements of }S\}.
\end{equation}
He singled out $S(2,10,26)$ because it is contained in the binary Golay lengths~\cite{Livinskyi2012}.  The multi-base theorem determines the order of growth of this function up to constants in the exponent.

\begin{theorem}[Covering by a fixed multiplicative monoid]\label{thm:ell-general}
Let $r\ge2$, let $b_1=2,b_2,\ldots,b_r>1$ be multiplicatively independent, and put
\[
S=S(b_1,\ldots,b_r)=\{b_1^{a_1}\cdots b_r^{a_r}:a_i\in\N_0\}.
\]
Then there is a constant $c>0$ such that
\begin{equation}\label{eq:ell-bounds}
\exp(c m\log m)
\le \ell_m(S)
\le \exp\bigl((r-1+o(1))m\log m\bigr).
\end{equation}
In particular,
\[
\log\ell_m(S)=\Theta(m\log m).
\]
\end{theorem}

\begin{proof}
\emph{Lower bound.}
Apply the upper-bound half of Theorem~\ref{thm:ksw} with digit set $\{0,1\}$.  The representability assumption holds because $2$ is one of the bases.  Hence every $n>2$ is a sum of at most
\[
K\frac{\log n}{\log\log n}
\]
distinct elements of $S$, for a fixed $K>0$.  Choose a sufficiently small constant $c>0$ and put
\[
N=\left\lfloor\exp(c m\log m)\right\rfloor.
\]
Since $\log x/\log\log x$ is eventually increasing, uniformly for $3\le n\le N$ we have
\[
K\frac{\log n}{\log\log n}
\le
K\frac{\log N}{\log\log N}
=(Kc+o(1))m.
\]
Taking, for instance, $c<1/(2K)$ shows that every $n\le N$ is a sum of at most $m$ elements of $S$ once $m$ is large.  Thus $\ell_m(S)>N$, proving the first inequality in \eqref{eq:ell-bounds}.

\smallskip
\noindent\emph{Upper bound.}
Let
\[
T(N):=|S\cap[1,N]|.
\]
If $b_1^{a_1}\cdots b_r^{a_r}\le N$, then
$a_i\le(\log N)/(\log b_i)$, so
\begin{equation}\label{eq:TN}
T(N)\le C_0(\log N)^r
\end{equation}
for a constant $C_0$ depending only on the bases.  Every representation of an integer at most $N$ as a sum of at most $m$ positive elements of $S$ uses only elements of $S\cap[1,N]$.  The number of multisets of size at most $m$ drawn from $T(N)$ available elements is at most
\[
\binom{T(N)+m}{m}.
\]
Fix $\eta>0$ and set
\[
N=\left\lceil\exp((r-1+\eta)m\log m)\right\rceil.
\]
By \eqref{eq:TN} and the standard estimate
$\binom{u+m}{m}\le (\ee(u+m)/m)^m$,
\begin{align*}
\log\binom{T(N)+m}{m}
&\le m\log\!\left(\frac{\ee(T(N)+m)}{m}\right)\\
&\le (r-1)m\log m+r m\log\log m+O(m).
\end{align*}
The right-hand side is strictly smaller than
\[
\log N=(r-1+\eta)m\log m+O(1)
\]
for all sufficiently large $m$.  Hence fewer than $N$ sums are available to cover the $N$ integers $1,\ldots,N$, so at least one is missing.  Therefore $\ell_m(S)\le N$.  Since $\eta>0$ was arbitrary, the second inequality in \eqref{eq:ell-bounds} follows.
\end{proof}

\begin{corollary}\label{cor:ell-21026}
For the Golay monoid $S(2,10,26)$,
\[
\ell_m(S(2,10,26))=\exp(\Theta(m\log m)),
\]
and more specifically
\[
\limsup_{m\to\infty}
\frac{\log \ell_m(S(2,10,26))}{m\log m}\le2.
\]
\end{corollary}

\begin{proof}
The bases $2,10,26$ are multiplicatively independent: from
\[
2^a10^b26^c=2^{a+b+c}5^b13^c=1
\]
with $a,b,c\in\mathbb Z$, unique factorization gives $b=c=0$ and then $a=0$.  Apply Theorem~\ref{thm:ell-general} with $r=3$.
\end{proof}

\begin{remark}
The same theorem applied to $S(2,3)$ gives
\[
\ell_m(S(2,3))=\exp(\Theta(m\log m)).
\]
Thus the very rapid initial values discussed by Livinskyi cannot have asymptotic growth of order $\exp(\Theta(m^2))$.  This is consistent with the modern sparse double-base representation theory that culminated in Theorem~\ref{thm:ksw}.
\end{remark}

\section{Gaps between Hadamard orders}\label{sec:gaps}
Let
\[
n_1<n_2<n_3<\cdots
\]
be the Hadamard orders and define, following Brent, Osborn and Smith~\cite{Brent2013},
\[
\gamma(x)=\max\{n_{i+1}-n_i:n_i\le x\}.
\]
A logarithmic estimate $\tau(q)\le\alpha\log_2q+O(1)$ gives a polynomial gap bound by comparing consecutive odd multipliers at a common power of two.  Theorem~\ref{thm:main} gives a subpolynomial version of the same idea.

\begin{theorem}[Quantitative subpolynomial Hadamard gaps]\label{thm:gaps}
As $x\to\infty$,
\begin{equation}\label{eq:gap}
\gamma(x)
\le
\exp\!\left(O\!\left(\frac{\log x}{\log\log x}\right)\right)
=x^{O(1/\log\log x)}
=x^{o(1)}.
\end{equation}
\end{theorem}

\begin{proof}
By Theorem~\ref{thm:main}, after enlarging constants to absorb finitely many odd $q$, there is a nondecreasing function
\[
F(q)=O\!\left(\frac{\log(q+3)}{\log\log(q+3)}\right)
\]
such that $2^tq$ is a Hadamard order for every odd $q$ and every integer $t\ge F(q)$.  The ``every $t$'' statement follows by tensoring with the Hadamard matrix of order $2$.

For large $x$, take
\[
T=\lceil F(2x)\rceil.
\]
Given any natural-number starting point $0\le y\le x$, choose the least positive odd integer $q$ for which
$2^Tq>y$.  Then $q\le y/2^T+2\le 2x$, so $T\ge F(q)$ and therefore $2^Tq$ is a Hadamard order.  By minimality of $q$,
\[
0<2^Tq-y\le2^{T+1}.
\]
Thus every point up to $x$ is followed by a Hadamard order at distance at most $2^{T+1}$, and hence
\[
\gamma(x)\le2^{T+1}
=\exp\!\left(O\!\left(\frac{\log x}{\log\log x}\right)\right).
\]
The two remaining forms in \eqref{eq:gap} are immediate.
\end{proof}

\begin{remark}[Relation to Smith's unpublished result]
Brent, Osborn and Smith reported in the 2013 version of their maximal-determinant paper that an unpublished manuscript of W.~D.~Smith proves
\[
\gamma(x)=O(x^{\eps})\qquad\text{for every }\eps>0,
\]
with potentially very large hidden constants~\cite{Brent2013}.  Consequently, the qualitative assertion $\gamma(x)=x^{o(1)}$ should not be presented as historically new.  The significance of Theorem~\ref{thm:gaps}, if the present transfer is not already in the literature, is that it supplies the explicit rate
$\exp(O(\log x/\log\log x))$ from published ingredients.
\end{remark}

\section{A consequence for approximate Hadamard matrices}\label{sec:approx}
Alexeev, Jasper and Mixon~\cite{Alexeev2025} study two relaxations of the Hadamard problem.  Let
\[
u(n)=\min_{M\in O(n)}\|M\|_{\max}
\]
be the smallest possible largest entry of an $n\times n$ orthogonal matrix, and let
\[
\kappa(n)=\min_{A\in\{\pm1\}^{n\times n}}
\frac{\sigma_{\max}(A)}{\sigma_{\min}(A)}
\]
be the best condition number of a sign matrix of order $n$.
They prove, using the then-current Hadamard-gap exponent $3/23$, that every fixed
$\alpha<17/92$ is admissible in
\[
\kappa(n)\le1+n^{-\alpha}
\]
for all sufficiently large $n$.  They further note that their method would allow every $\alpha<1/4$ under the Hadamard conjecture.  Theorem~\ref{thm:gaps} reaches this limiting exponent without assuming the conjecture.

\begin{corollary}[Flat orthogonal matrices]\label{cor:flat}
For every fixed $0<\eps<1/2$, all sufficiently large $n$ satisfy
\[
u(n)\le \frac{1}{\sqrt n-n^{\eps}}.
\]
\end{corollary}

\begin{proof}
The restriction $0<\eps<1/2$ ensures that the displayed denominator is positive for all sufficiently large $n$.  Theorem~\ref{thm:gaps} implies $\gamma(n)<n^{\eps}$ for all sufficiently large $n$.  Hence there is a Hadamard order
\[
m\in[n,n+n^{\eps}).
\]
Put $k=m-n<n^{\eps}$.  For large $n$, $k<\sqrt m$.  Lemmas~3 and~4 of Alexeev--Jasper--Mixon, applied to an $n\times n$ submatrix of a Hadamard matrix of order $m$, produce an orthogonal matrix with maximum entry at most
\[
\frac{1}{\sqrt m-k}
\le \frac{1}{\sqrt n-n^{\eps}}.
\]
\end{proof}

\begin{corollary}[Approximate Hadamard condition numbers]\label{cor:kappa}
For every $\delta>0$, all sufficiently large $n$ satisfy
\begin{equation}\label{eq:kappa}
\kappa(n)\le1+n^{-1/4+\delta}.
\end{equation}
\end{corollary}

\begin{proof}
The random-rounding argument in Alexeev--Jasper--Mixon shows that a bound of the form in Corollary~\ref{cor:flat} with exponent $\eps$ yields
\[
\kappa(n)\le1+n^{-\alpha}
\]
for every fixed $\alpha<1/4-\eps/2$ and all sufficiently large $n$~\cite{Alexeev2025}.  Given $\delta>0$, choose $\eps>0$ so small that $\eps<2\delta$ and then take
$\alpha=1/4-\delta$.  This gives \eqref{eq:kappa}.
\end{proof}

\begin{remark}
The exponent in Corollary~\ref{cor:kappa} is the limiting value of the Alexeev--Jasper--Mixon gap-based method even if the full Hadamard conjecture is assumed.  Their paper emphasizes that other explicit constructions suggest that the true optimal decay for $\kappa(n)-1$ may be faster; Corollary~\ref{cor:kappa} concerns only what follows from improving the Hadamard-order gaps.
\end{remark}

\section{Effectivity, limitations, and further problems}
The argument above is short because two powerful existing results fit together unusually well.  That same modularity makes its present limitations clear.

\paragraph{Effective constants.}
Theorem~\ref{thm:ksw} is asymptotic, and Krenn--Suppakitpaisarn--Wagner explicitly note that the finite preprocessing threshold required by their greedy proof can be very large.  Thus Theorem~\ref{thm:main} does not presently provide a competitive numerical construction for a specified moderate $p$.  A useful next result would be an effective version for the particular bases $(2,10)$, with explicit constants and a realistic threshold.  Because the multi-base proof is driven by Diophantine approximation between logarithms of the bases, continued fractions and explicit irrationality measures for $\log 10/\log 2$ are natural tools.

\paragraph{Sharper covering constants.}
For $S(2,10,26)$ we proved only
\[
\exp(c m\log m)\le\ell_m(S)\le\exp((2+o(1))m\log m).
\]
Determining the optimal constants in
\[
\frac{\log\ell_m(S)}{m\log m}
\]
would refine Livinskyi's additive program and might yield useful finite decompositions even without changing the $\Theta(m\log m)$ order.

\paragraph{Beyond a fixed Golay monoid.}
The lower-bound half of the multi-base theorem indicates that $\Theta(\log n/\log\log n)$ is intrinsic to fixed finite collections of multiplicatively independent bases in this representation model.  To beat Theorem~\ref{thm:main} by the same broad strategy one should therefore exploit more than a fixed finite multiplicative monoid: for example, additional primitive Golay lengths, complementary sets of bounded cardinality at many new lengths, or additive constructions that do not reduce to fixed-base power-products.

\paragraph{Reducing the transfer overhead.}
Craigen's theorem converts $r$ complementary pairs into a factor $4^{r+1}$ in the Hadamard order.  Any construction replacing this by $2^{o(r)}$, while retaining sufficiently flexible lengths, would immediately improve the exponent in Theorem~\ref{thm:main}.  Conversely, an obstruction to such a transfer would help identify the real bottleneck.

\paragraph{Bibliographic verification.}
The central transfer
\[
\text{sparse multi-base expansion}
\Longrightarrow
\text{sparse sum of Golay lengths}
\Longrightarrow
\text{Hadamard matrix}
\]
is elementary once Theorems~\ref{thm:ksw} and~\ref{thm:craigen-transfer} are placed side by side.  That makes a thorough priority check especially important.  The present targeted search found no source making this connection, but a submission should not rely on search-engine absence as evidence of novelty.

\section*{Acknowledgement of research status}
This manuscript is a research draft intended for mathematical verification and prior-art review.  In particular, statements advertised as contributions should be independently checked before public claims of novelty are made.


\section*{Formalization scope and proof alignment}
The accompanying Lean project proves the final unconditional sparse
representation, Hadamard existence, least-exponent, gap, flatness and
conditioning results using mathlib and project-owned sources. The proof
imports contain no holes or additional axioms. Its 357 principal-declaration
axiom audit, including 204 quantitative logarithm lemmas, uses only
Lean's three standard axioms.

The formal sparse argument proves that an irrational ratio of logarithms of
natural integers at least two is not Liouville. Integer auxiliary coefficients,
block binomial polynomials, controlled denominators, stable complex
interpolation and twenty finite extrapolation stages establish the separation
input; rotation meshes and a greedy argument then give sparse representations.
This route avoids separately formalizing the arbitrary-digit theorem cited
as background in the manuscript. All additional proof code is independently
written from mathlib.

The file \texttt{Challenge.lean} independently specifies six statements using
only mathlib. Its intentional placeholder proofs are excluded from every
proof import. The completed proofs are exported by \texttt{Solution.lean};
local comparison and verification evidence is recorded in
\path{submission/verification.json}. Local checking is not registry
acceptance, independent human review, or a determination of novelty.

The construction checks binary Golay seeds of lengths $1,2,10,26$, Turyn
multiplication, and the Craigen consequence with order $4^{r+1}(2L+1)$.
Gap predicates use natural starting points, and flatness requires
$0<\varepsilon<1/2$. Mathematical citations are background references rather
than imported proof code. See \texttt{DEPENDENCIES.md} and
\path{paper/ALIGNMENT.md} for precise scope.

\begin{thebibliography}{99}

\bibitem{Alexeev2025}
B. Alexeev, J. Jasper, and D. G. Mixon,
\emph{Asymptotically optimal approximate Hadamard matrices},
arXiv:2511.14653, 2025.

\bibitem{Brent2013}
R. P. Brent, J.-A. H. Osborn, and W. D. Smith,
\emph{Lower bounds on maximal determinants of $\{\pm1\}$ matrices via the probabilistic method},
arXiv:1211.3248v3, 2013.


\bibitem{Craigen1995}
R. Craigen,
\emph{Signed groups, sequences, and the asymptotic existence of Hadamard matrices},
Journal of Combinatorial Theory, Series A \textbf{71} (1995), 241--254.
\href{https://doi.org/10.1016/0097-3165(95)90002-0}{doi:10.1016/0097-3165(95)90002-0}.

\bibitem{deLauney2009}
W. de Launey,
\emph{On the asymptotic existence of Hadamard matrices},
Journal of Combinatorial Theory, Series A \textbf{116} (2009), 1002--1008.
\href{https://doi.org/10.1016/j.jcta.2009.01.001}{doi:10.1016/j.jcta.2009.01.001}.

\bibitem{DuJiang2024}
C. Du and Y. Jiang,
\emph{Golay complementary sequences of arbitrary length and asymptotic existence of Hadamard matrices},
arXiv:2401.15381, 2024.

\bibitem{Golay1961}
M. J. E. Golay,
\emph{Complementary series},
IRE Transactions on Information Theory \textbf{7} (1961), 82--87.
\href{https://doi.org/10.1109/TIT.1961.1057620}{doi:10.1109/TIT.1961.1057620}.

\bibitem{JedwabParker2009}
J. Jedwab and M. G. Parker,
\emph{A construction of binary Golay sequence pairs from odd-length Barker sequences},
Journal of Combinatorial Designs \textbf{17} (2009), 478--491.
\href{https://doi.org/10.1002/jcd.20222}{doi:10.1002/jcd.20222}.

\bibitem{Krenn2020}
D. Krenn, V. Suppakitpaisarn, and S. Wagner,
\emph{On the minimal Hamming weight of a multi-base representation},
Journal of Number Theory \textbf{208} (2020), 168--179.
\href{https://doi.org/10.1016/j.jnt.2019.07.023}{doi:10.1016/j.jnt.2019.07.023}.

\bibitem{Livinskyi2012}
I. Livinskyi,
\emph{Asymptotic Existence of Hadamard Matrices},
M.Sc. thesis, University of Manitoba, 2012.

\bibitem{Turyn1974}
R. J. Turyn,
\emph{Hadamard matrices, Baumert--Hall units, four-symbol sequences, pulse compression, and surface wave encodings},
Journal of Combinatorial Theory, Series A \textbf{16} (1974), 313--333.
\href{https://doi.org/10.1016/0097-3165(74)90056-9}{doi:10.1016/0097-3165(74)90056-9}.

\bibitem{Wallis1976}
J. S. Wallis,
\emph{On the existence of Hadamard matrices},
Journal of Combinatorial Theory, Series A \textbf{21} (1976), 188--195.
\href{https://doi.org/10.1016/0097-3165(76)90062-5}{doi:10.1016/0097-3165(76)90062-5}.

\end{thebibliography}

\end{document}
